\documentclass[11pt,a4paper]{article}
\usepackage[utf8]{inputenc}
\usepackage[T1]{fontenc}
\usepackage[english]{babel}
\usepackage{amsmath,amssymb,amsthm}
\usepackage{forest}
\usepackage{booktabs}
\usepackage{enumitem}
\usepackage{longtable}
\usepackage[margin=2.6cm]{geometry}
\usepackage[colorlinks=true,linkcolor=blue!60!black,citecolor=green!40!black]{hyperref}

\title{Multirate infinitesimal step methods with polynomial coupling:\\
a B-series theory of order conditions and its subclasses}
\author{Oswald Knoth}
\date{\today}

\newtheorem{theorem}{Theorem}[section]
\newtheorem{proposition}[theorem]{Proposition}
\newtheorem{lemma}[theorem]{Lemma}
\newtheorem{corollary}[theorem]{Corollary}
\theoremstyle{definition}
\newtheorem{definition}[theorem]{Definition}
\newtheorem{example}[theorem]{Example}
\theoremstyle{remark}
\newtheorem{remark}[theorem]{Remark}

\newcommand{\R}{\mathbb{R}}
\newcommand{\T}{\mathcal{T}}
\newcommand{\dd}{\,\mathrm{d}}

% bicoloured trees (root at the bottom, root colour irrelevant)
\forestset{
  b/.style={circle,draw=black,fill=black,inner sep=1.3pt},
  w/.style={circle,draw=black,fill=white,inner sep=1.3pt},
  bitree/.style={for tree={grow=90,s sep=3pt,l=5pt,l sep=0pt,
                 parent anchor=center,child anchor=center,edge={thin}},
                 baseline}
}
\newcommand{\btree}[1]{\begin{forest}bitree #1\end{forest}}
\newcommand{\Tb}{\btree{[,b]}}
\newcommand{\Tbb}{\btree{[,b [,b]]}}
\newcommand{\Tbw}{\btree{[,b [,w]]}}
\newcommand{\Tbbb}{\btree{[,b [,b][,b]]}}
\newcommand{\Tbbw}{\btree{[,b [,b][,w]]}}
\newcommand{\Tbww}{\btree{[,b [,w][,w]]}}
\newcommand{\Tbbbl}{\btree{[,b [,b [,b]]]}}
\newcommand{\Tbbwl}{\btree{[,b [,b [,w]]]}}
\newcommand{\Tbwbl}{\btree{[,b [,w [,b]]]}}
\newcommand{\Tbwwl}{\btree{[,b [,w [,w]]]}}
\newcommand{\Tfourbbw}{\btree{[,b [,b][,b][,w]]}}
\newcommand{\Tfourbww}{\btree{[,b [,b][,w][,w]]}}
\newcommand{\Tfourbbl}{\btree{[,b [,b [,b]][,w]]}}
\newcommand{\Tfourbbwl}{\btree{[,b [,b [,b][,w]]]}}

\begin{document}
\maketitle

\begin{abstract}
Multirate infinitesimal step (MIS) methods integrate an ordinary differential equation $y'=F(y,y)$ whose first argument
varies slowly and whose second argument carries the fast dynamics: the first argument is frozen at stage values of a
Runge--Kutta-like outer scheme, and the second argument is integrated exactly in auxiliary (``infinitesimal'') initial
value problems. We study MIS methods in which both the coupling to the frozen right-hand sides and the coupling to the
increments of previous stages are polynomials in the inner time. Using B-series over trees with bicoloured non-root
vertices, whose coefficients are polynomials in the inner time, we derive the order conditions in a Runge--Kutta-like
form, give them explicitly up to order four, and provide an algorithm in exact arithmetic for arbitrary order. The
general theory specialises to the known theories of several method classes: Runge--Kutta methods (the base method),
multirate infinitesimal GARK methods, classical and generalised split-explicit MIS methods, commutator-free Lie group
methods including their low-storage variant, and explicit exponential Runge--Kutta methods, whose non-stiff conditions are the conditions over a subset of the trees and whose stiff order
conditions become algebraic conditions on the coupling polynomials. For methods whose inner problems start from a
previous stage we prove that order three requires only the order-three conditions of the base method and a single
coupling condition, and that the purely fast conditions hold automatically. As examples, the exponential
Runge--Kutta method ETDRK4-B of Krogstad, read as an MIS method, has order four for the general problem $y'=F(y,y)$, and
for the low-storage commutator-free methods of Bazavov the order-three question reduces to a single condition, which we
observe to follow from the order of the underlying 2N-storage Runge--Kutta method.
\end{abstract}

\tableofcontents

% =====================================================================
\section{Introduction}
\label{sec:intro}

Many evolution problems contain processes on widely separated time scales. In atmospheric models for weather and
climate prediction, and in low-Mach-number compressible flow, the solution of the compressible Euler equations is
dominated by slow advection and physical processes, while sound waves (and, at high resolution, gravity waves) are fast
but carry little energy. An explicit integrator would have to resolve the fast waves. Two remedies are in use:
(semi-)implicit treatment of the fast terms, e.g.\ horizontally explicit/vertically implicit (HEVI) and IMEX schemes
\cite{GiraldoKellyConstantinescu2013,BoscarinoFilbetRusso2016}, and \emph{split-explicit} schemes, in which the slow
tendencies are evaluated with a large step and held fixed while the fast terms are integrated with many small steps
\cite{KlempWilhelmson1978,WickerSkamarock2002,KlempSkamarockDudhia2007,Baldauf2010}.

Split-explicit schemes do not split the right-hand side into a sum of a slow and a fast term. They freeze the slowly
varying \emph{quantities}: in the compressible Euler equations, for example, the advecting velocity and the transported
potential temperature are frozen at a stage value, while mass fluxes and pressure are integrated with small steps. This
is formalised by writing the semi-discrete equation $y'=\mathcal F(y)$ as
\begin{equation}\label{eq:G}
y'=F(y,y),\qquad F:\R^N\times\R^N\to\R^N,\qquad F(y,y)=\mathcal F(y), \tag{$\mathcal G$}
\end{equation}
where the first argument carries the slow (frozen) and the second argument the fast (integrated) dependence. A method
evaluates $F(Y_j,Z)$ with a stage value $Y_j$ and the solution $Z$ of an inner initial value problem. The
\emph{multirate infinitesimal step} (MIS) methods \cite{KnothWolke1998,WenschKnothGalant2009,KnothWensch2014}
provide the mathematical framework for such schemes: in the analysis the inner problems are solved exactly, which
gives order conditions for the outer method. In applications the inner problems are solved by small steps of a cheap
explicit or partly implicit method.

In this paper the coupling of the inner problems to the frozen right-hand sides and to the increments of previous stages
is a \emph{polynomial in the inner time}. This class contains the classical MIS and split-explicit methods (constant
coupling), the multirate infinitesimal GARK (MRI-GARK) methods of Sandu \cite{Sandu2019} and related multirate
exponential methods \cite{LuanChinomonaReynolds2020,ChinomonaReynolds2021}, commutator-free Lie group methods
\cite{CelledoniMarthinsenOwren2003,Owren2006} and their low-storage form \cite{Bazavov2022}, and a class of explicit
exponential Runge--Kutta methods \cite{HochbruckOstermann2005,HochbruckOstermann2010}. Our aim is a single order theory for the
whole class, from which the theories of these subclasses follow by restricting the set of trees or the coefficients.

\paragraph{Contributions.}
\begin{itemize}
\item B-series for MIS methods with polynomial coupling over trees with bicoloured non-root vertices
  (Theorem~\ref{thm:recursion}), and the order conditions in Runge--Kutta-like form (Theorem~\ref{thm:general}), explicit
  up to order four.
\item The restriction to the problem classes affine, additive and linear in the fast argument, with the numbers of
  conditions (Table~\ref{tab:treecount}).
\item For methods whose inner problems start from a previous stage: the purely fast conditions hold automatically, and
  order three is equivalent to order three of the base method plus one coupling condition (Theorem~\ref{thm:exactfast}),
  without the uniform-speed (internal consistency) assumption.
\item The reduction to known theories: Runge--Kutta methods (base method), constant-coupling MIS methods, MRI-GARK,
  commutator-free Lie group methods and their low-storage variant, and explicit exponential Runge--Kutta methods
  including their stiff order conditions (Section~\ref{sec:subclasses}).
\end{itemize}

\paragraph{Outline.}
Section~\ref{sec:problem} introduces the problem classes, Section~\ref{sec:methods} the methods and their special cases,
Section~\ref{sec:trees} the trees and B-series, Section~\ref{sec:oc} the order conditions, and
Section~\ref{sec:subclasses} the specialisation to subclasses with an existing theory. Section~\ref{sec:algorithm}
describes the computation of the conditions, and Section~\ref{sec:conclusions} lists open questions. The order-four
conditions are given in Appendix~\ref{app:order4}.

% =====================================================================
\section{Problem classes}
\label{sec:problem}

We consider \eqref{eq:G} with $F$ sufficiently smooth. Throughout, the index $y$ ($z$) at a derivative denotes
differentiation with respect to the first (second) argument. Important special cases are
\begin{align}
 &\text{affine in the fast argument:} && F(y,z)=f(y)+B(y)\,z, \tag{$\mathcal B$}\label{eq:B}\\
 &\text{additive:} && F(y,z)=f(y)+g(z), \tag{$\mathcal A$}\label{eq:A}\\
 &\text{constant linear fast part:} && F(y,z)=f(y)+Lz. \tag{$\mathcal L$}\label{eq:L}
\end{align}
In \eqref{eq:B} the fast operator depends on the frozen state, as in the split-explicit treatment of the Euler equations;
if $B$ depends linearly on a set of frozen coefficient fields, a linear combination $\sum_jp_{ij}F(Y_j,z)$ needs only one
application of $B$ to $z$, which makes polynomial coupling cheap. Every system \eqref{eq:G} can be written as a
partitioned system by doubling the unknowns, $u'=F(u,v)$, $v'=F(u,v)$, $u(0)=v(0)$, so partitioned B-series (P-series)
apply \cite{HairerNorsettWanner1993}; since both components are equal, the colour of the root plays no role.

% =====================================================================
\section{MIS methods with polynomial coupling}
\label{sec:methods}

Throughout, $\tau\in[0,1]$ denotes the normalised time of the inner problems.

\begin{definition}[MIS method with polynomial coupling]\label{def:mis}
Let $s\ge1$, and for $2\le i\le s+1$, $1\le j<i$ let $\alpha_{ij}\in\R$ and polynomials
\[
 p_{ij}(\tau)=\sum_{k=0}^{K}a_{ijk}\,\tau^k,\qquad g_{ij}(\tau)=\sum_{k=0}^{K}\gamma_{ijk}\,\tau^k
\]
be given. One step $y_n\mapsto y_{n+1}$ is defined by $Y_1=y_n$ and, for $i=2,\dots,s+1$,
\begin{align}
 Z_i(0) &= y_n+\sum_{j=1}^{i-1}\alpha_{ij}\,(Y_j-y_n), \label{eq:mis1}\\
 Z_i'(\tau) &= \sum_{j=1}^{i-1}g_{ij}(\tau)\,(Y_j-y_n) + h\sum_{j=1}^{i-1}p_{ij}(\tau)\,F\big(Y_j,Z_i(\tau)\big),
   \qquad \tau\in[0,1], \label{eq:mis2}\\
 Y_i &= Z_i(1), \label{eq:mis3}
\end{align}
and $y_{n+1}=Y_{s+1}$. $K$ is the \emph{coupling degree}; $K=0$ is \emph{constant coupling}.
\end{definition}

The fast part is integrated with the speed $q_i(\tau):=\sum_jp_{ij}(\tau)$ over the fast time
$Q_i(1)$, $Q_i(\tau):=\int_0^\tau q_i$. We call
\begin{equation}\label{eq:uniform}
 q_i(\tau)\equiv d_i\ \text{(constant)},\qquad\text{i.e.}\quad \sum_j a_{ij0}=d_i,\quad \sum_ja_{ijk}=0\ (k\ge1),
\end{equation}
\emph{uniform fast speed}. It does not make the inner problem autonomous, since the frozen arguments still enter with
time-dependent weights.

\subsection{Stage times and the base method}

Let $P_{ij}(\tau)=\int_0^\tau p_{ij}$, $G_{ij}(\tau)=\int_0^\tau g_{ij}$. The \emph{stage times} are defined by the method
applied to $y'=1$:
\begin{equation}\label{eq:ci}
 c_1=0,\qquad c_i=\sum_{j<i}\big(\alpha_{ij}+G_{ij}(1)\big)\,c_j+\sum_{j<i}P_{ij}(1).
\end{equation}
If $F$ does not depend on its second argument, the inner problems are quadratures and
$Y_i-y_n=\sum_j(\alpha_{ij}+G_{ij}(1))(Y_j-y_n)+h\sum_jP_{ij}(1)F(Y_j)$. With the strictly lower triangular
$(s+1)\times(s+1)$ matrices $\bar A=(\alpha_{ij}+G_{ij}(1))$, $\bar P=(P_{ij}(1))$ this is an explicit Runge--Kutta
method, the \emph{base method}, with
\begin{equation}\label{eq:base}
 \begin{pmatrix}A^{RK}\\ b^{RK\,T}\end{pmatrix}=(I-\bar A)^{-1}\bar P .
\end{equation}
The shape of $p_{ij}$ and $g_{ij}$ enters only through the dependence on the fast argument.

\subsection{Special cases}
\label{sec:specialcases}

\paragraph{Classical MIS and generalised split-explicit methods.}
In \cite{WenschKnothGalant2009} the inner problem is posed on $[0,d_ih]$ with constant coefficients $\beta_{ij}$,
$d_i=\sum_j\beta_{ij}$; after rescaling to $\tau\in[0,1]$ this is Definition~\ref{def:mis} with $K=0$,
$p_{ij}\equiv\beta_{ij}$ and uniform fast speed. The generalised split-explicit methods of \cite{KnothWensch2014} add
constant $g_{ij}\equiv\gamma_{ij}$.

\paragraph{Split-explicit Runge--Kutta methods.}
The RK3 scheme of Wicker and Skamarock \cite{WickerSkamarock2002} integrates every stage from $y_n$ with the slow
tendency of the previous stage: $\alpha=0$, $g\equiv0$, $p_{i,i-1}\equiv c_i$ with $c=(0,\tfrac13,\tfrac12,1)$, all
other $p_{ij}=0$.

\paragraph{Telescopic methods, MRI-GARK.}
If $\alpha_{i,i-1}=1$, all other $\alpha_{ij}=0$ and $g\equiv0$, then $Z_i(0)=Y_{i-1}$, and the stages are inner
integrations over consecutive fast time intervals of length $c_i-c_{i-1}$. With polynomial $p_{ij}$ and additive $F$
these are the explicit MRI-GARK methods of Sandu \cite{Sandu2019}; uniform fast speed corresponds to their internal
consistency conditions. Related classes are the multirate exponential Runge--Kutta methods
\cite{LuanChinomonaReynolds2020} and the IMEX-MRI-GARK methods \cite{ChinomonaReynolds2021}.

\paragraph{Commutator-free Lie group methods.}
A commutator-free method \cite{CelledoniMarthinsenOwren2003,Owren2006} computes its stages as compositions of flows of
frozen vector fields,
\begin{equation}\label{eq:cf}
 Y_i=\exp\Big(h\sum_j\beta^{(J)}_{ij}F_j\Big)\circ\dots\circ\exp\Big(h\sum_j\beta^{(1)}_{ij}F_j\Big)\,y_n,\qquad
 F_j=F(Y_j,\cdot),
\end{equation}
where $\exp(hX)y$ is the flow of the vector field $X$ over the time $h$. Each exponential is the solution of an inner
problem \eqref{eq:mis2} with constant coupling $p_{ij}\equiv\beta^{(l)}_{ij}$, $g\equiv0$, that starts from the result of the
previous exponential. Hence \eqref{eq:cf} is an MIS method with constant coupling whose stages are the exponentials; the
intermediate exponentials are auxiliary stages, which are not used as frozen arguments (their columns of $p$ vanish).
For the frame representation $F(y,z)=\sum_kf_k(y)E_k(z)$ of a problem on a homogeneous space these are the
commutator-free Lie group methods; for general $F$ they are MIS methods for \eqref{eq:G}.

\paragraph{Exponential Runge--Kutta methods.}
For \eqref{eq:L} the inner problems are solved by the variation-of-constants formula. With
\[
 \varphi_{k+1}(z)=\int_0^1e^{(1-\tau)z}\frac{\tau^k}{k!}\dd\tau,\qquad k\ge0,
\]
a method with $\alpha=0$, $g\equiv0$ and uniform fast speed $q_i\equiv c_i$ gives
\begin{equation}\label{eq:exprk}
 Y_i=e^{c_ihL}y_n+h\sum_{j<i}a_{ij}(hL)f(Y_j),\qquad
 a_{ij}(z)=\int_0^1e^{(1-\tau)c_iz}p_{ij}(\tau)\dd\tau=\sum_{k=0}^{K}k!\,a_{ijk}\,\varphi_{k+1}(c_iz).
\end{equation}

\begin{proposition}\label{prop:exprk}
An explicit exponential Runge--Kutta method $Y_i=e^{c_ihL}y_n+h\sum_ja_{ij}(hL)f(Y_j)$, $y_{n+1}=Y_{s+1}$,
$c_{s+1}=1$, is an MIS method with polynomial coupling, $\alpha=0$ and $g\equiv0$, if and only if every $a_{ij}(z)$ is a
linear combination of $\varphi_k(c_iz)$, $k\ge1$, and $\sum_ja_{ij}(z)=c_i\varphi_1(c_iz)$.
\end{proposition}

\begin{proof}
By \eqref{eq:exprk} and the linear independence of $\varphi_1(cz),\varphi_2(cz),\dots$ for fixed $c\ne0$. The condition
$\sum_ja_{ij}=c_i\varphi_1(c_iz)$ is equivalent to $q_i\equiv c_i$.
\end{proof}

Such methods applied to \eqref{eq:G} are a nonlinear generalisation of exponential Runge--Kutta methods, in which the
exponential is replaced by the flow of the fast part with frozen slow argument.

\begin{example}\label{ex:exp}
\begin{enumerate}[label=(\alph*)]
\item Exponential Euler, $y_{n+1}=e^{hL}y_n+h\varphi_1(hL)f(y_n)$: $s=1$, $p_{21}=1$.
\item ETD2RK of Cox and Matthews \cite{CoxMatthews2002}, $c_2=1$, $a_{21}=\varphi_1$, $b_1=\varphi_1-\varphi_2$,
  $b_2=\varphi_2$: $p_{21}=1$, $p_{31}=1-\tau$, $p_{32}=\tau$; the slow tendency is linearly interpolated in time between
  the two stages.
\item ETDRK4-B of Krogstad \cite{Krogstad2005}, $c=(0,\tfrac12,\tfrac12,1)$, with $\varphi_{k,i}=\varphi_k(c_iz)$:
  $a_{21}=\tfrac12\varphi_{1,2}$, $a_{31}=\tfrac12\varphi_{1,3}-\varphi_{2,3}$, $a_{32}=\varphi_{2,3}$,
  $a_{41}=\varphi_{1,4}-2\varphi_{2,4}$, $a_{43}=2\varphi_{2,4}$, $b_1=\varphi_1-3\varphi_2+4\varphi_3$,
  $b_2=b_3=2\varphi_2-4\varphi_3$, $b_4=4\varphi_3-\varphi_2$. In MIS form
  \begin{gather*}
   p_{21}=\tfrac12,\quad p_{31}=\tfrac12-\tau,\ p_{32}=\tau,\quad p_{41}=1-2\tau,\ p_{43}=2\tau,\\
   p_{51}=1-3\tau+2\tau^2,\quad p_{52}=p_{53}=2\tau-2\tau^2,\quad p_{54}=-\tau+2\tau^2,
  \end{gather*}
  all other $p_{ij}=0$ (the factor $k!$ of \eqref{eq:exprk}: $\varphi_{k+1}\leftrightarrow\tau^k/k!$).
\end{enumerate}
\end{example}

Without uniform fast speed the closed form remains available,
\[
 Y_i=e^{Q_i(1)hL}y_n+h\sum_j\int_0^1e^{(Q_i(1)-Q_i(\tau))hL}p_{ij}(\tau)\dd\tau\,f(Y_j),
\]
but the weights are no longer
combinations of $\varphi$-functions; with $\alpha\ne0$ products $e^{d_iz}\varphi_k(c_jz)$ appear.

\paragraph{Methods not contained.}
Integrating-factor (Lawson) methods \cite{Lawson1967} use $e^{(c_i-c_j)hL}f(Y_j)$: the slow tendency is transported
by the fast flow. In MIS form this would need an impulsive forcing at $\tau=c_j/c_i$, which is not polynomial.
Semi-implicit schemes of the type \cite{BoscarinoFilbetRusso2016}, which also freeze the first argument and treat the
second implicitly, are MIS-type methods in which the inner problem is replaced by a single implicit step.

\subsection{Methods that restart from a previous stage}

A subclass that contains most of the methods above consists of methods whose inner problems start at a previous stage
and contain no increments of previous stages:
\begin{equation}\label{eq:restart}
 g_{ij}\equiv0,\qquad Z_i(0)=Y_{k(i)}\ \text{for some }k(i)<i
\end{equation}
(row $i$ of $\alpha$ is the unit vector $e_{k(i)}^T$, or zero if $k(i)=1$). It contains the classical MIS methods with
$Z_i(0)=y_n$, the telescopic methods, the commutator-free methods \eqref{eq:cf}, and the exponential Runge--Kutta methods of Proposition~\ref{prop:exprk}; the
coupling polynomials are arbitrary, uniform fast speed is not required. Following the chain
$s+1\to k(s+1)\to k(k(s+1))\to\dots\to1$ we call its stages the \emph{chain} of the method.

% =====================================================================
\section{Bicoloured trees and B-series}
\label{sec:trees}

\subsection{Trees and elementary differentials}

Let $\T$ be the set of rooted trees whose non-root vertices are coloured black or white. The colour of a vertex states
which argument of the $F$ at its parent is differentiated: black the first (frozen) argument, white the second
(integrated) one. Every tree is written $t=[t_1,\dots,t_m]$, where the subtrees $t_k$ carry the colour
$\kappa(t_k)\in\{\mathrm b,\mathrm w\}$ of their roots. The elementary differentials of \eqref{eq:G} are
\[
 \mathbf F(\bullet)=F(y,y),\qquad
 \mathbf F([t_1,\dots,t_m])=\partial_{\kappa(t_1)}\cdots\partial_{\kappa(t_m)}F(y,y)\big[\mathbf F(t_1),\dots,\mathbf F(t_m)\big],
\]
with $\partial_{\mathrm b}=\partial_y$, $\partial_{\mathrm w}=\partial_z$. Order $|t|$, symmetry $\sigma(t)$ (automorphisms
preserving colours) and density $\gamma(t)=|t|\prod_k\gamma(t_k)$ are defined as usual. Since $y'=F(y,y)$, the exact
solution is
\[
 y(t_n+h)=y_n+\sum_{t\in\T}\frac{h^{|t|}}{\sigma(t)\gamma(t)}\mathbf F(t)(y_n),
\]
independent of the colours. For the problem classes of Section~\ref{sec:problem} the set of trees shrinks:
\begin{itemize}
\item \eqref{eq:B}: $\partial_z^2F=0$, so a vertex has at most one white child;
\item \eqref{eq:A}: $\partial_y\partial_zF=0$, so all children of a vertex have the same colour;
\item \eqref{eq:L}: in addition $\partial_z^2F=0$, so a vertex has either only black children or exactly one, white,
  child: the white vertices form chains ($L$-chains), the trees of exponential integrators;
\item trees with only black vertices give the base method, trees with only white non-root vertices the purely fast
  problem $y'=g(y)$.
\end{itemize}

\begin{table}[ht]
\centering
\caption{Number of trees per order (non-root vertices bicoloured, root colour irrelevant), i.e.\ the number of order
conditions at each order, for the problem classes. Computed by enumeration.}
\label{tab:treecount}
\begin{tabular}{lrrrrrrr}
\toprule
order & 1 & 2 & 3 & 4 & 5 & 6 & 7\\\midrule
general \eqref{eq:G} & 1 & 2 & 7 & 26 & 107 & 458 & 2058\\
affine in $z$ \eqref{eq:B} & 1 & 2 & 6 & 20 & 73 & 278 & 1106\\
additive \eqref{eq:A} & 1 & 2 & 6 & 18 & 60 & 204 & 734\\
constant linear \eqref{eq:L} & 1 & 2 & 5 & 13 & 37 & 108 & 332\\
classical RK (black only) & 1 & 1 & 2 & 4 & 9 & 20 & 48\\
\bottomrule
\end{tabular}
\end{table}

\subsection{B-series with polynomial coefficients}

Write the stage values and the inner solutions as B-series,
\[
 Y_i=B(\eta_i,y_n),\qquad Z_i(\tau)=B\big(\zeta_i(\tau),y_n\big),\qquad
 B(a,y)=a(\emptyset)y+\sum_{t\in\T}\frac{h^{|t|}a(t)}{\sigma(t)}\mathbf F(t)(y),
\]
with $\eta_i(\emptyset)=\zeta_i(\tau;\emptyset)=1$. The coefficients $\zeta_i(\tau;t)$ are polynomials in $\tau$.

\begin{theorem}\label{thm:recursion}
The coefficients of the method of Definition~\ref{def:mis} satisfy, for $t=[t_1,\dots,t_m]$, $2\le i\le s+1$,
$\eta_1(t)=0$,
\begin{align*}
 \zeta_i(0;t) &= \sum_{j<i}\alpha_{ij}\,\eta_j(t),\\
 \frac{\dd}{\dd\tau}\zeta_i(\tau;t) &= \sum_{j<i}g_{ij}(\tau)\,\eta_j(t)
   + \sum_{j<i}p_{ij}(\tau)\prod_{k=1}^m\begin{cases}\eta_j(t_k), & \kappa(t_k)=\mathrm b,\\
                  \zeta_i(\tau;t_k), & \kappa(t_k)=\mathrm w,\end{cases}\\
 \eta_i(t) &= \zeta_i(1;t).
\end{align*}
The method has order $p$ if and only if $\eta_{s+1}(t)=1/\gamma(t)$ for all $t\in\T$ with $|t|\le p$ (for generic $F$).
\end{theorem}

\begin{proof}
Insert the B-series of $Y_j$ and $Z_i(\tau)$ into $hF(Y_j,Z_i(\tau))$ and expand in a Taylor series in both arguments.
The coefficient of $\mathbf F([t_1,\dots,t_m])$ is the product of the coefficients of the subtrees, each taken from the
series of the argument that is differentiated; this is the composition of P-series \cite{HairerNorsettWanner1993}.
Integrating in $\tau$ and using \eqref{eq:mis1}, \eqref{eq:mis3} gives the recursion. The equivalence with order $p$
follows from the independence of the elementary differentials for generic $F$.
\end{proof}

The degree of $\zeta_i(\cdot;t)$ is at most $|t|(K+1)$; all integrals are exact.

% =====================================================================
\section{Order conditions}
\label{sec:oc}

\subsection{General form}

Let $q_i=\sum_jp_{ij}$, $Q_i=\sum_jP_{ij}$, $\bar A$, $\bar P$ as in \eqref{eq:base}, and
\[
 \Lambda=(I-\bar A)^{-1},\qquad \hat b^T=e_{s+1}^T\Lambda,\qquad \mathfrak A=\Lambda\bar P .
\]
The stage times are $c=\Lambda(Q_i(1))_i$, and the \emph{inner time} of stage $i$ is the polynomial
\[
 z_i(\tau):=\zeta_i(\tau;\bullet)=\sum_{j<i}\big(\alpha_{ij}+G_{ij}(\tau)\big)c_j+Q_i(\tau),\qquad z_i(1)=c_i .
\]
$\mathfrak A$ is the coefficient matrix of the base method \eqref{eq:base}: its rows $2,\dots,s$ are the stage
coefficients, its last row is $b^{RK\,T}$.

\begin{theorem}[general form of the order conditions]\label{thm:general}
For $t=[t_1,\dots,t_m]$ let
\begin{equation}\label{eq:dit}
 d_i(t)=\sum_{j<i}\int_0^1p_{ij}(\sigma)\prod_{k:\,\kappa(t_k)=\mathrm b}\eta_j(t_k)\prod_{k:\,\kappa(t_k)=\mathrm w}\zeta_i(\sigma;t_k)\dd\sigma .
\end{equation}
Then $\eta(t)=\Lambda\,d(t)$ and
$\zeta_i(\tau;t)=\sum_{j<i}(\alpha_{ij}+G_{ij}(\tau))\eta_j(t)+\sum_{j<i}\int_0^\tau p_{ij}(\sigma)\prod\cdots\dd\sigma$ with the
products of \eqref{eq:dit}. The method has order $p$ if and only if
\begin{equation}\label{eq:oc}
 \hat b^Td(t)=\sum_{i=2}^{s+1}\hat b_i\,d_i(t)=\frac1{\gamma(t)}\qquad\text{for all }t\in\T,\ |t|\le p .
\end{equation}
\end{theorem}

\begin{proof}
Integrating the recursion of Theorem~\ref{thm:recursion} from $0$ to $1$ gives $\eta_i(t)=\sum_j\bar A_{ij}\eta_j(t)+d_i(t)$,
i.e.\ $(I-\bar A)\eta(t)=d(t)$. The formula for $\zeta_i$ is the integral from $0$ to $\tau$.
\end{proof}

The structure is that of the Runge--Kutta conditions $b^T\Phi(t)=1/\gamma(t)$: $\hat b$ plays the role of the weights,
$d_i(t)$ integrates the product of the subtree coefficients against the coupling polynomials, black subtrees contribute
stage values $\eta_j$ of the coupled stage $j$, and white subtrees the inner coefficients $\zeta_i(\sigma;\cdot)$ of the
current stage $i$ along the inner time. Every condition is a polynomial equation in $\alpha_{ij}$, $a_{ijk}$,
$\gamma_{ijk}$.

\subsection{Explicit conditions up to order three}

Table~\ref{tab:oc3} lists $d_i(t)$ for the ten trees of order $\le3$; the conditions are $\hat b^Td(t)=1/\gamma(t)$. The
auxiliary quantities are
\begin{align*}
 \eta(\Tbb)&=\mathfrak Ac, &
 \eta(\Tbw)&=\Lambda\Big(\int_0^1q_iz_i\Big)_i,\\
 \zeta_i(\sigma;\Tbb)&=\sum_j\big(\alpha_{ij}+G_{ij}(\sigma)\big)\eta_j(\Tbb)+\sum_jP_{ij}(\sigma)c_j, &
 \zeta_i(\sigma;\Tbw)&=\sum_j\big(\alpha_{ij}+G_{ij}(\sigma)\big)\eta_j(\Tbw)+\int_0^\sigma q_iz_i .
\end{align*}

\begin{table}[ht]
\centering
\caption{Order conditions $\sum_i\hat b_id_i(t)=1/\gamma(t)$ up to order three (trees drawn with the root at the
bottom). Column ``class'' refers to Theorem~\ref{thm:exactfast}: B base method condition, F automatically satisfied,
C coupling condition, -- implied.}
\label{tab:oc3}
\renewcommand{\arraystretch}{1.6}
\begin{tabular}{rcllc}
\toprule
order & $t$ & $d_i(t)$ & $1/\gamma(t)$ & class\\\midrule
1 & \Tb & $Q_i(1)$ & $1$ & B\\\midrule
2 & \Tbb & $\sum_jP_{ij}(1)\,c_j$ & $\frac12$ & B\\
  & \Tbw & $\int_0^1q_i(\sigma)\,z_i(\sigma)\dd\sigma$ & $\frac12$ & F\\\midrule
3 & \Tbbb & $\sum_jP_{ij}(1)\,c_j^2$ & $\frac13$ & B\\
  & \Tbbw & $\sum_jc_j\int_0^1p_{ij}(\sigma)\,z_i(\sigma)\dd\sigma$ & $\frac13$ & C\\
  & \Tbww & $\int_0^1q_i(\sigma)\,z_i(\sigma)^2\dd\sigma$ & $\frac13$ & F\\
  & \Tbbbl & $\sum_jP_{ij}(1)\,\eta_j(\Tbb)$ & $\frac16$ & B\\
  & \Tbbwl & $\sum_jP_{ij}(1)\,\eta_j(\Tbw)$ & $\frac16$ & --\\
  & \Tbwbl & $\int_0^1q_i(\sigma)\,\zeta_i(\sigma;\Tbb)\dd\sigma$ & $\frac16$ & --\\
  & \Tbwwl & $\int_0^1q_i(\sigma)\,\zeta_i(\sigma;\Tbw)\dd\sigma$ & $\frac16$ & F\\
\bottomrule
\end{tabular}
\end{table}

The conditions of order four (26 trees) are given in Appendix~\ref{app:order4}.

\subsection{The base method}

\begin{proposition}\label{prop:base}
For trees with only black vertices, $\eta(t)$ is the Runge--Kutta elementary weight of the base method
$\mathfrak A=\Lambda\bar P$. In particular, order $p$ of the MIS method implies order $p$ of the base method.
\end{proposition}

\begin{proof}
For black trees no $\zeta_i(\sigma;\cdot)$ appears in \eqref{eq:dit}, so $d_i(t)=\sum_jP_{ij}(1)\prod_k\eta_j(t_k)$ and
$\eta(t)=\Lambda d(t)$ is the recursion of the elementary weights of the Runge--Kutta method $\mathfrak A$.
\end{proof}

\subsection{Methods that restart from a previous stage}

Under \eqref{eq:restart}, $\Lambda$ has a simple structure: $\eta_i=\eta_{k(i)}+d_i$, hence $\hat b_i=1$ for the stages
of the chain and $\hat b_i=0$ otherwise, and the inner time is $z_i(\sigma)=c_{k(i)}+Q_i(\sigma)$.

\begin{theorem}\label{thm:exactfast}
Assume \eqref{eq:restart} and $c_{s+1}=1$.
\begin{enumerate}[label=(\alph*)]
\item For every tree $u$ whose non-root vertices are all white, $\eta_i(u)=c_i^{|u|}/\gamma(u)$ and
  $\zeta_i(\sigma;u)=z_i(\sigma)^{|u|}/\gamma(u)$. In particular all purely fast conditions are satisfied.
\item For all stages
  \[
   \eta_i(\Tbbwl)=\tfrac12\,\eta_i(\Tbbb),\qquad \eta_i(\Tbbw)+\eta_i(\Tbwbl)=c_i\,\eta_i(\Tbb).
  \]
\item The method has order three if and only if the base method has order three and the coupling condition
  \begin{equation}\label{eq:coupling3}
   \sum_{i\in\text{chain}}\ \sum_{j<i}c_j\int_0^1p_{ij}(\sigma)\big(c_{k(i)}+Q_i(\sigma)\big)\dd\sigma=\frac13
  \end{equation}
  holds.
\end{enumerate}
\end{theorem}

\begin{proof}
(a) For $F(y,z)=g(z)$ the inner problem is $Z_i'=hq_i(\tau)g(Z_i)$, $Z_i(0)=Y_{k(i)}$, the flow of $g$ reparametrised in
time. Hence $Z_i(\sigma)$ is the exact flow over the time $h(c_{k(i)}+Q_i(\sigma))=hz_i(\sigma)$, provided $Y_{k(i)}$ is the
exact flow over $hc_{k(i)}$. Induction over $i$ gives the claim, since the coefficients of the exact flow over the time
$hc$ are $c^{|u|}/\gamma(u)$.

(b) By (a), $\eta_j(\Tbw)=c_j^2/2$, so $d_i(\Tbbwl)=\tfrac12d_i(\Tbbb)$. For the second identity, by (a) and
integration by parts with $Q_i'=q_i$, $P_{ij}'=p_{ij}$, and $k=k(i)$,
\begin{align*}
 d_i(\Tbbw)+d_i(\Tbwbl)
 &=\sum_jc_j\int_0^1p_{ij}\big(c_{k}+Q_i\big)+\int_0^1q_i\Big(\eta_{k}(\Tbb)+\sum_jP_{ij}c_j\Big)\\
 &=c_k\sum_jP_{ij}(1)c_j+\sum_jc_j\big[P_{ij}Q_i\big]_0^1+Q_i(1)\eta_k(\Tbb)
  =c_i\eta_i(\Tbb)-c_k\eta_k(\Tbb),
\end{align*}
using $c_i=c_k+Q_i(1)$ and $\eta_i(\Tbb)=\eta_k(\Tbb)+\sum_jP_{ij}(1)c_j$. Since $\eta_i=\eta_k+d_i$, induction over $i$
gives the identity.

(c) By (a) the conditions for \Tbw, \Tbww, \Tbwwl\ hold. By (b), $\eta_{s+1}(\Tbbwl)=\tfrac16$ if and only if
$\eta_{s+1}(\Tbbb)=\tfrac13$, and, given $\eta_{s+1}(\Tbb)=\tfrac12$, $\eta_{s+1}(\Tbwbl)=\tfrac16$ if and only if
$\eta_{s+1}(\Tbbw)=\tfrac13$, which is \eqref{eq:coupling3}. The remaining conditions \Tb, \Tbb, \Tbbb, \Tbbbl\ are those
of the base method.
\end{proof}

\begin{proposition}[computational]\label{prop:count}
Under \eqref{eq:restart} the numbers of coupling conditions that are independent of the base method conditions and of
the conditions of lower order are
\begin{center}
\begin{tabular}{lcccc}
\toprule
order & 2 & 3 & 4 & 5\\\midrule
all trees & 2 & 7 & 26 & 107\\
base method (black trees) & 1 & 2 & 4 & 9\\
coupling conditions, class \eqref{eq:restart} & 0 & 1 & 4 & 16\\
coupling conditions, general $\alpha$, $g_{ij}$ & 1 & 5 & 22 & 89\\
\bottomrule
\end{tabular}
\end{center}
At order four a possible choice of the four coupling conditions is (with $\hat b_i=1$ on the chain, $k=k(i)$,
$z_i(\sigma)=c_{k}+Q_i(\sigma)$)
\begin{align*}
 &\Tfourbbw: && \sum_i\hat b_i\sum_jc_j^2\int_0^1p_{ij}z_i=\tfrac14, &
 &\Tfourbww: && \sum_i\hat b_i\sum_jc_j\int_0^1p_{ij}z_i^2=\tfrac14,\\
 &\Tfourbbl: && \sum_i\hat b_i\sum_j\eta_j(\Tbb)\int_0^1p_{ij}z_i=\tfrac18, &
 &\Tfourbbwl: && \sum_i\hat b_i\sum_jP_{ij}(1)\,\eta_j(\Tbbw)=\tfrac1{12},
\end{align*}
with $\eta_j(\Tbbw)=\sum_{l\in\text{chain}(j)}\sum_rc_r\int_0^1p_{lr}z_l$. All other mixed trees of order four are implied
(class -- in Table~\ref{tab:oc4}).
\end{proposition}

The counts were obtained as ranks of the stage-wise coefficient vectors of many random methods, computed exactly modulo
a large prime, with the products of lower-order coefficients included as admissible relations. For orders two and three
they agree with Theorem~\ref{thm:exactfast}; a proof of the order-four statement in the style of
Theorem~\ref{thm:exactfast}(b) is open.

% =====================================================================
\section{Subclasses with an existing order theory}
\label{sec:subclasses}

The general theory reduces to the known theories of several subclasses in two ways: by restricting the set of trees to
the problem class (Section~\ref{sec:trees}) and by restricting the coefficients of the method.

\subsection{Runge--Kutta methods}
If $F$ does not depend on $z$, only black trees remain, and by Proposition~\ref{prop:base} the conditions
\eqref{eq:oc} are the Runge--Kutta conditions of the base method \eqref{eq:base}. Conversely, every explicit
Runge--Kutta method $(A,b)$ is the base method of the MIS method with $\alpha=0$, $g\equiv0$ and constant
$p_{ij}\equiv a_{ij}$, and of its telescopic form with $p_{ij}\equiv a_{ij}-a_{i-1,j}$, $\alpha_{i,i-1}=1$.

\subsection{Constant coupling: classical MIS and split-explicit methods}
For $K=0$ all polynomials are constants, $z_i(\sigma)$ is linear in $\sigma$, and every integral in \eqref{eq:dit} is
explicit. The conditions of Table~\ref{tab:oc3} then reduce to those of
\cite{WenschKnothGalant2009,KnothWensch2014} for constant coupling and are extended to all orders by
Theorem~\ref{thm:general}; constant-coupling conditions were extended in \cite{BauerKnoth2021}. The split-explicit RK3
method of Wicker and Skamarock is of order two for \eqref{eq:G} and \eqref{eq:A}: its base method has $b=(0,0,1)$,
$c=(0,\tfrac13,\tfrac12)$, so the condition $b^Tc^2=\tfrac13$ of \Tbbb\ fails ($b^Tc^2=\tfrac14$). Its base method has linear order three,
but with a fast part a coupling condition of order three fails as well (the residual of \Tbwbl\ is $\tfrac1{12}$), so
it is of order two also for \eqref{eq:L}.

\subsection{Additive splitting and MRI-GARK}
For \eqref{eq:A} the conditions are those of the trees in which all children of a vertex have the same colour
(Table~\ref{tab:treecount}). For telescopic methods with polynomial coupling these are MRI-GARK methods, for which
Sandu \cite{Sandu2019} derived the order conditions via GARK theory \cite{GuntherSandu2016}: the conditions of the slow
base method plus coupling conditions, under the internal consistency (uniform fast speed) assumption. Telescopic methods
satisfy \eqref{eq:restart}, so Theorem~\ref{thm:exactfast} applies: the purely fast conditions hold for exact inner
solutions, and order three needs only the base method and \eqref{eq:coupling3}, with or without uniform fast speed. The
uniform fast speed remains relevant for the inner integrator, since $q_i<0$ would integrate the fast part backwards in
time. As a check of the reduction, the recursion of Theorem~\ref{thm:recursion} gives order three for MRI-GARK-ERK33 and
order four for MRI-GARK-ERK45a \cite{Sandu2019}. In contrast to the GARK-based derivation, the general theory also gives
the conditions for non-additive $F$, i.e.\ for frozen slow arguments inside the fast part.

\subsection{Commutator-free Lie group methods}
\label{sec:cf}

By Section~\ref{sec:specialcases}, a commutator-free method \eqref{eq:cf} is an MIS method of the restart class
\eqref{eq:restart} with constant coupling. Theorem~\ref{thm:general} applied to it gives its order conditions for
arbitrary frozen vector fields $F(Y_j,\cdot)$ on $\R^N$, i.e.\ for \eqref{eq:G}. Owren \cite{Owren2006} derived the order
conditions of commutator-free Lie group methods in a free Lie algebra. Since the Lie algebra of all vector fields on
$\R^N$ satisfies no relations beyond those of a free Lie algebra in the generic case, we expect both sets of conditions to
characterise the same methods; a formal proof is open, and in all examples below the orders agree. By
Theorem~\ref{thm:exactfast}, applied to the MIS form:
\begin{itemize}
\item the purely fast conditions hold automatically (the exponentials are exact flows);
\item a commutator-free method has order three if and only if its base method --- the Runge--Kutta method obtained by
  replacing every exponential by the Euler step $y\mapsto y+h\sum_j\beta^{(l)}_{ij}F_j$ --- has order three and the single
  coupling condition \eqref{eq:coupling3} holds.
\end{itemize}

\begin{example}
CF4 of Celledoni, Marthinsen and Owren \cite{CelledoniMarthinsenOwren2003},
\begin{gather*}
 Y_2=\exp(\tfrac12hF_1)y_n,\quad Y_3=\exp(\tfrac12hF_2)y_n,\quad Y_4=\exp\big(h(F_3-\tfrac12F_1)\big)\exp(\tfrac12hF_1)y_n,\\
 y_{n+1}=\exp\big(\tfrac h{12}(-F_1+2F_2+2F_3+3F_4)\big)\exp\big(\tfrac h{12}(3F_1+2F_2+2F_3-F_4)\big)y_n,
\end{gather*}
is in MIS form a method with $s=6$ stages (two auxiliary). The recursion in exact arithmetic gives order four for
\eqref{eq:G}, \eqref{eq:B}, \eqref{eq:A} and \eqref{eq:L}, its order as a Lie group method.
\end{example}

\paragraph{Low-storage commutator-free methods.}
Bazavov \cite{Bazavov2022} proposed commutator-free methods built on a Williamson 2N-storage Runge--Kutta method
\cite{Williamson1980}, $\Delta Y_i=A_i\Delta Y_{i-1}+hF(Y_{i-1})$, $Y_i=Y_{i-1}+B_i\Delta Y_i$ ($A_1=0$):
\begin{equation}\label{eq:lscf}
 Z_i=A_iZ_{i-1}+hF_{i-1},\qquad Y_i=\exp(B_iZ_i)\,Y_{i-1},\qquad i=1,\dots,s,
\end{equation}
which needs storage for two quantities and one exponential per stage. In the notation of Definition~\ref{def:mis} (stages
shifted by one) this is the telescopic method
\[
 \alpha_{i+1,i}=1,\qquad g\equiv0,\qquad p_{i+1,j}\equiv B_i\prod_{k=j+1}^{i}A_k\quad(j\le i),
\]
with the 2N-storage Runge--Kutta method as base method. The three-stage third-order method of L\"uscher for the Wilson
flow \cite{Luescher2010} is of this form ($A=(0,-\tfrac{17}{32},-\tfrac{32}{27})$, $B=(\tfrac14,\tfrac89,\tfrac34)$).
Bazavov proved that three-stage methods \eqref{eq:lscf} of classical order three are Lie group methods of order three,
and conjectured that this holds for more stages and higher orders \cite{Bazavov2022,BazavovChuna2021}. With
Theorem~\ref{thm:exactfast} the order-three part of the conjecture reduces to a single statement:

\begin{corollary}\label{cor:lscf}
A method \eqref{eq:lscf} has order three for \eqref{eq:G} if and only if its 2N-storage Runge--Kutta method has order
three and \eqref{eq:coupling3} holds, which here reads
\[
 \sum_{i=1}^{s}\sum_{j\le i}p_{i+1,j}\,c_j\,\frac{c_i+c_{i+1}}2=\frac13,
\]
with $c_j$ the stage times of the MIS form (the frozen argument $F_j$ is taken at time $c_j$).
\end{corollary}

For random methods \eqref{eq:lscf} with $s=4$ and $s=5$ stages whose Runge--Kutta method has order three (40 methods each,
obtained by a Newton iteration on the four classical conditions), all seven bicoloured conditions of order three hold to
round-off; that is, \eqref{eq:coupling3} follows from the classical conditions in every case. A proof for general $s$ is
open. At order four, the recursion gives order four for \eqref{eq:G} for the 2N-storage method RK(5,4) of Carpenter and
Kennedy \cite{CarpenterKennedy1994} in the form \eqref{eq:lscf}, and for a random six-stage method whose Runge--Kutta
method has order four, in agreement with the conjecture.

\subsection{Explicit exponential Runge--Kutta methods}
\label{sec:exponential}

For \eqref{eq:L} only the trees in which the white vertices form $L$-chains remain. Their elementary differentials are
$\mathbf F(t)=L^{l_1}f^{(m)}(\dots)$, products of powers of $L$ and derivatives of $f$, i.e.\ the elementary
differentials of exponential integrators in their non-stiff expansion in powers of $hL$. For methods with $\alpha=0$,
$g\equiv0$ and uniform fast speed $q_i\equiv c_i$, Proposition~\ref{prop:exprk} identifies the MIS method with an
exponential Runge--Kutta method; Theorem~\ref{thm:general} restricted to these trees therefore gives the classical
(non-stiff) order conditions of explicit exponential Runge--Kutta methods.

\begin{example}
For ETD2RK (Example~\ref{ex:exp}(b)) the recursion gives order two in all classes; its residuals
$\eta_3(t)-1/\gamma(t)$ at order three are
\begin{center}
\begin{tabular}{cccccccc}
\toprule
$t$ & \Tbbb & \Tbbw & \Tbww & \Tbbbl & \Tbbwl & \Tbwbl & \Tbwwl\\\midrule
residual & $\frac16$ & $0$ & $0$ & $-\frac16$ & $\frac1{12}$ & $0$ & $0$\\
\bottomrule
\end{tabular}
\end{center}
in agreement with Theorem~\ref{thm:exactfast}: the purely fast conditions and \eqref{eq:coupling3} hold, the order is
limited by the base method (Heun's method), and the residual of \Tbbwl\ is half that of \Tbbb. For ETDRK4-B
(Example~\ref{ex:exp}(c)) the recursion in exact rational arithmetic gives order four for \eqref{eq:L} --- the classical
order of the method --- and order four also for \eqref{eq:A}, \eqref{eq:B} and the general problem \eqref{eq:G}. Read
as an MIS method with quadratic coupling, ETDRK4-B is therefore a fourth-order method for $y'=F(y,y)$ with a nonlinear,
state-dependent fast part, if the inner problems are solved exactly.
\end{example}

\paragraph{Stiff order conditions.}
The classical conditions assume that $h\|\partial_zF\|$ is small. Hochbruck and Ostermann
\cite{HochbruckOstermann2005,HochbruckOstermann2010} derived \emph{stiff order conditions} for exponential Runge--Kutta
methods, which have to hold uniformly in $hL$. With $b_i=a_{s+1,i}$ and
\[
 \psi_j(z)=\sum_ib_i(z)\frac{c_i^{j-1}}{(j-1)!}-\varphi_j(z),\qquad
 \psi_{j,i}(z)=\sum_ka_{ik}(z)\frac{c_k^{j-1}}{(j-1)!}-c_i^j\varphi_j(c_iz),
\]
the conditions up to order three read $\psi_1(hL)=0$; $\psi_2(hL)=0$, $\psi_{1,i}(hL)=0$; $\psi_3(hL)=0$,
$\sum_ib_i(hL)J\psi_{2,i}(hL)=0$ for arbitrary bounded operators $J$, and analogous conditions at order four
\cite{HochbruckOstermann2005}. For MIS methods with polynomial coupling they become algebraic conditions on the
coefficients:

\begin{proposition}\label{prop:stiff}
Let $\alpha=0$, $g\equiv0$, $q_i\equiv c_i$, and define the polynomials
\[
 r_i^{[j]}(\tau)=\sum_{k<i}p_{ik}(\tau)\frac{c_k^{j-1}}{(j-1)!}-c_i^j\frac{\tau^{j-1}}{(j-1)!} .
\]
\begin{enumerate}[label=(\alph*)]
\item $\psi_{j,i}\equiv0$ if and only if $r_i^{[j]}\equiv0$, and $\psi_j\equiv0$ if and only if $r_{s+1}^{[j]}\equiv0$.
  In particular $\psi_{1,i}\equiv0$ is the uniform fast speed condition.
\item A sufficient condition for $\sum_ib_i(z)J\psi_{2,i}(z)=0$ for all $J$ is that for every value $c$ of the stage times
  \[
   \sum_{i:\,c_i=c}p_{s+1,i}(\tau)\,r_i^{[2]}(\sigma)\equiv0\qquad\text{as a polynomial in }(\tau,\sigma).
  \]
\end{enumerate}
\end{proposition}

\begin{proof}
(a) $\psi_{j,i}(z)=\int_0^1e^{(1-\tau)c_iz}r_i^{[j]}(\tau)\dd\tau$, and the map from polynomials to these functions is
injective for $c_i\ne0$. (b) $b_i(z_1)\psi_{2,i}(z_2)=\iint e^{(1-\tau)z_1}e^{(1-\sigma)c_iz_2}p_{s+1,i}(\tau)r_i^{[2]}(\sigma)\dd\tau\dd\sigma$;
since $J$ is arbitrary, the two arguments have to be treated as independent.
\end{proof}

For ETD2RK, $r_3^{[2]}=(1-\tau)\cdot0+\tau\cdot1-\tau=0$, so it has stiff order two. Analogous bivariate and trivariate
polynomial identities, grouped by pairs of stage times, are sufficient for the order-four conditions.

\subsection{Affine fast part}
For \eqref{eq:B}, trees with a vertex with two or more white children drop out: 20 instead of 26 conditions at order four
and 73 instead of 107 at order five (Table~\ref{tab:treecount}). A method constructed for \eqref{eq:B} only has to satisfy
these. The class contains \eqref{eq:L} and is the structure of frozen fast operators that depend on the slow state.

% =====================================================================
\section{Computing the order conditions}
\label{sec:algorithm}

Theorem~\ref{thm:recursion} gives an algorithm that is independent of the special structure and works in exact rational
or symbolic arithmetic. The trees are processed in increasing order, so that all subtree coefficients are available.

\begin{center}
\begin{minipage}{0.9\linewidth}
\begin{tabbing}
xx\=xx\=xx\=xx\=xx\=\kill
\textbf{for} each tree $t=[t_1,\dots,t_m]$, ordered by $|t|$\\
\>\textbf{for} $i=2,\dots,s+1$\\
\>\>$\phi(\tau)\gets\sum_{j<i}g_{ij}(\tau)\eta_j(t)$\\
\>\>\textbf{for} $j=1,\dots,i-1$\\
\>\>\>$r(\tau)\gets1$\\
\>\>\>\textbf{for} $k=1,\dots,m$: \ $r\gets r\cdot\zeta_i(\cdot;t_k)$ if $t_k$ is white (polynomial product), else $r\gets r\cdot\eta_j(t_k)$\\
\>\>\>$\phi\gets\phi+p_{ij}\,r$\\
\>\>$\zeta_i(\tau;t)\gets\sum_{j<i}\alpha_{ij}\eta_j(t)+\int_0^\tau\phi(\tilde\tau)\dd\tilde\tau$\\
\>\>$\eta_i(t)\gets\zeta_i(1;t)$\\
\textbf{residual:} $\eta_{s+1}(t)-1/\gamma(t)$
\end{tabbing}
\end{minipage}
\end{center}

The algorithm needs only the information which argument is differentiated; restricting the enumeration of trees to the
admissible child sets of Section~\ref{sec:trees} gives the conditions of the subclasses. With symbolic coefficients the
residuals are polynomial equations in $\alpha_{ij}$, $a_{ijk}$, $\gamma_{ijk}$, the starting point for the construction
of new methods; with dual numbers they provide exact gradients for numerical searches. Independent implementations exist
in Fortran, in Python, in Julia (the module \texttt{MIS} of the package \texttt{MultiRate}), and, for constant $g_{ij}$,
in \texttt{BSeries.jl} (\texttt{MultirateInfinitesimalSplitMethod}) \cite{KetchesonRanocha2023}. They reproduce
Table~\ref{tab:treecount}, Tables~\ref{tab:oc3} and~\ref{tab:oc4}, Theorem~\ref{thm:exactfast} (checked up to order five),
and the orders of the methods cited in Section~\ref{sec:subclasses}.

% =====================================================================
\section{Conclusions and open questions}
\label{sec:conclusions}

The B-series over trees with bicoloured non-root vertices and polynomial coefficients give a single order theory for MIS
methods with polynomial coupling. Restricting the trees to the problem class and the coefficients to a subclass yields
the order conditions of Runge--Kutta methods, constant-coupling MIS methods, MRI-GARK methods, commutator-free Lie group
methods and explicit exponential Runge--Kutta methods; for methods that restart from a previous stage the purely fast conditions are automatic and order
three needs a single coupling condition. Open questions are:
\begin{enumerate}
\item \emph{Simplifying assumptions.} Is there an analogue of the stage-order assumption $C(q)$ that reduces the
  conditions to those of the base method plus few coupling conditions? A natural candidate is the polynomial identity
  $\sum_{j<i}p_{ij}(\tau)c_j^{k-1}=c_i^k\tau^{k-1}$, $k=1,\dots,q$ (for $\alpha=0$, $g\equiv0$), which is the strong form of
  the stiff conditions in Proposition~\ref{prop:stiff}; as for Runge--Kutta methods it cannot hold for $i=2$, $k\ge2$.
\item \emph{Order four in the restart class.} A proof of Proposition~\ref{prop:count}.
\item \emph{Stiff order.} Necessity of the conditions in Proposition~\ref{prop:stiff}(b); conditions for skew-adjoint
  (hyperbolic) fast operators, where no parabolic smoothing is available; a stiff theory for nonlinear fast parts, e.g.\
  by B-series whose coefficients are functions of $h\partial_zF$.
\item \emph{Inexact inner solutions.} If the inner problems are solved by $m$ steps of a Runge--Kutta method of order $q$,
  the integration $\int_0^\tau\phi$ in the algorithm is replaced by that method, and one expects
  $\eta^{(m)}_{s+1}(t)=\eta_{s+1}(t)+\sum_{k\ge q}e_k(t)m^{-k}$. For order $p$ with every $m$ the inner method has to
  reproduce $\eta_{s+1}(t)$ exactly for $|t|\le p$. We conjecture that for $K=0$, $g\equiv0$ an inner method of order
  $q\ge p$ suffices, while $K>0$ may require a higher inner (quadrature) order.
\item \emph{Low-storage commutator-free methods.} A proof that \eqref{eq:coupling3} follows from the classical order-three
  conditions for every 2N-storage method (Corollary~\ref{cor:lscf}), and the analogous statement at order four, which
  would prove Bazavov's conjecture; a formal identification of Owren's free-Lie-algebra conditions with the conditions
  of Theorem~\ref{thm:general} for the frame representation.
\item \emph{Further subclasses.} Exponential methods with mixed arguments $\varphi_k(c_jz)$ in stage $i$ and exponential Rosenbrock
  methods \cite{HochbruckOstermannSchweitzer2009} are MIS methods only for special coefficients, possibly with
  $\alpha\ne0$.
\end{enumerate}

% =====================================================================
\appendix
\section{Order conditions of order four}
\label{app:order4}

Conditions $\sum_i\hat b_id_i(t)=1/\gamma(t)$ for the 26 trees of order four, general method of
Definition~\ref{def:mis}. $\eta_j(u)$ and $\zeta_i(\sigma;u)$ for $|u|\le3$ are given by Theorem~\ref{thm:general} and
Table~\ref{tab:oc3}. Column ``class'' refers to the subclass \eqref{eq:restart}: B base method, F automatically
satisfied, C1--C4 the coupling conditions of Proposition~\ref{prop:count}, -- implied by lower order, base and C1--C4.

{\renewcommand{\arraystretch}{1.5}
\begin{longtable}{rcllc}
\caption{Order-four conditions.}\label{tab:oc4}\\
\toprule
No. & $t$ & $d_i(t)$ & $1/\gamma(t)$ & class\\\midrule
\endfirsthead
\toprule
No. & $t$ & $d_i(t)$ & $1/\gamma(t)$ & class\\\midrule
\endhead
\bottomrule
\endfoot
1 & \btree{[,b [,b [,b [,b]]]]} & $\sum_j P_{ij}(1)\,\eta_j(\Tbbbl)$ & $\frac1{24}$ & B\\[3pt]
2 & \btree{[,b [,b [,b [,w]]]]} & $\sum_j P_{ij}(1)\,\eta_j(\Tbbwl)$ & $\frac1{24}$ & --\\[3pt]
3 & \btree{[,b [,b [,b][,b]]]} & $\sum_j P_{ij}(1)\,\eta_j(\Tbbb)$ & $\frac1{12}$ & B\\[3pt]
4 & \btree{[,b [,b [,b][,w]]]} & $\sum_j P_{ij}(1)\,\eta_j(\Tbbw)$ & $\frac1{12}$ & C4\\[3pt]
5 & \btree{[,b [,b [,b]][,w]]} & $\sum_j \eta_j(\Tbb)\int_0^1 p_{ij}(\sigma)\,z_i(\sigma)\dd\sigma$ & $\frac1{8}$ & C3\\[3pt]
6 & \btree{[,b [,b [,w [,b]]]]} & $\sum_j P_{ij}(1)\,\eta_j(\Tbwbl)$ & $\frac1{24}$ & --\\[3pt]
7 & \btree{[,b [,b [,w [,w]]]]} & $\sum_j P_{ij}(1)\,\eta_j(\Tbwwl)$ & $\frac1{24}$ & --\\[3pt]
8 & \btree{[,b [,b [,w][,w]]]} & $\sum_j P_{ij}(1)\,\eta_j(\Tbww)$ & $\frac1{12}$ & --\\[3pt]
9 & \btree{[,b [,b [,w]][,w]]} & $\sum_j \eta_j(\Tbw)\int_0^1 p_{ij}(\sigma)\,z_i(\sigma)\dd\sigma$ & $\frac1{8}$ & --\\[3pt]
10 & \btree{[,b [,b][,b [,b]]]} & $\sum_j P_{ij}(1)\,c_j\eta_j(\Tbb)$ & $\frac1{8}$ & B\\[3pt]
11 & \btree{[,b [,b][,b [,w]]]} & $\sum_j P_{ij}(1)\,c_j\eta_j(\Tbw)$ & $\frac1{8}$ & --\\[3pt]
12 & \btree{[,b [,b][,b][,b]]} & $\sum_j P_{ij}(1)\,c_j^{3}$ & $\frac1{4}$ & B\\[3pt]
13 & \btree{[,b [,b][,b][,w]]} & $\sum_j c_j^{2}\int_0^1 p_{ij}(\sigma)\,z_i(\sigma)\dd\sigma$ & $\frac1{4}$ & C1\\[3pt]
14 & \btree{[,b [,b][,w [,b]]]} & $\sum_j c_j\int_0^1 p_{ij}(\sigma)\,\zeta_i(\sigma;\Tbb)\dd\sigma$ & $\frac1{8}$ & --\\[3pt]
15 & \btree{[,b [,b][,w [,w]]]} & $\sum_j c_j\int_0^1 p_{ij}(\sigma)\,\zeta_i(\sigma;\Tbw)\dd\sigma$ & $\frac1{8}$ & --\\[3pt]
16 & \btree{[,b [,b][,w][,w]]} & $\sum_j c_j\int_0^1 p_{ij}(\sigma)\,\big(z_i(\sigma)\big)^{2}\dd\sigma$ & $\frac1{4}$ & C2\\[3pt]
17 & \btree{[,b [,w [,b [,b]]]]} & $\int_0^1 q_i(\sigma)\,\zeta_i(\sigma;\Tbbbl)\dd\sigma$ & $\frac1{24}$ & --\\[3pt]
18 & \btree{[,b [,w [,b [,w]]]]} & $\int_0^1 q_i(\sigma)\,\zeta_i(\sigma;\Tbbwl)\dd\sigma$ & $\frac1{24}$ & --\\[3pt]
19 & \btree{[,b [,w [,b][,b]]]} & $\int_0^1 q_i(\sigma)\,\zeta_i(\sigma;\Tbbb)\dd\sigma$ & $\frac1{12}$ & --\\[3pt]
20 & \btree{[,b [,w [,b][,w]]]} & $\int_0^1 q_i(\sigma)\,\zeta_i(\sigma;\Tbbw)\dd\sigma$ & $\frac1{12}$ & --\\[3pt]
21 & \btree{[,b [,w [,w [,b]]]]} & $\int_0^1 q_i(\sigma)\,\zeta_i(\sigma;\Tbwbl)\dd\sigma$ & $\frac1{24}$ & --\\[3pt]
22 & \btree{[,b [,w [,w [,w]]]]} & $\int_0^1 q_i(\sigma)\,\zeta_i(\sigma;\Tbwwl)\dd\sigma$ & $\frac1{24}$ & F\\[3pt]
23 & \btree{[,b [,w [,w][,w]]]} & $\int_0^1 q_i(\sigma)\,\zeta_i(\sigma;\Tbww)\dd\sigma$ & $\frac1{12}$ & F\\[3pt]
24 & \btree{[,b [,w][,w [,b]]]} & $\int_0^1 q_i(\sigma)\,z_i(\sigma)\zeta_i(\sigma;\Tbb)\dd\sigma$ & $\frac1{8}$ & --\\[3pt]
25 & \btree{[,b [,w][,w [,w]]]} & $\int_0^1 q_i(\sigma)\,z_i(\sigma)\zeta_i(\sigma;\Tbw)\dd\sigma$ & $\frac1{8}$ & F\\[3pt]
26 & \btree{[,b [,w][,w][,w]]} & $\int_0^1 q_i(\sigma)\,\big(z_i(\sigma)\big)^{3}\dd\sigma$ & $\frac1{4}$ & F\\[3pt]
\end{longtable}}

\begin{thebibliography}{99}
\small
\bibitem{Bazavov2022} A.~Bazavov, Commutator-free Lie group methods with minimum storage requirements and reuse of
 exponentials, \emph{BIT Numer. Math.} (2022), arXiv:2007.04225.
\bibitem{BazavovChuna2021} A.~Bazavov, T.~Chuna, Efficient integration of gradient flow in lattice gauge theory and
 properties of low-storage commutator-free Lie group methods, arXiv:2101.05320 (2021).
\bibitem{Baldauf2010} M.~Baldauf, Linear stability analysis of Runge--Kutta-based partial time-splitting schemes for the
 Euler equations, \emph{Mon. Wea. Rev.} 138 (2010) 4475--4496.
\bibitem{BauerKnoth2021} S.~Bauer, O.~Knoth, Extended multirate infinitesimal step methods: derivation of order
 conditions, \emph{J. Comput. Appl. Math.} 387 (2021) 112541.
\bibitem{BoscarinoFilbetRusso2016} S.~Boscarino, F.~Filbet, G.~Russo, High order semi-implicit schemes for time
 dependent partial differential equations, \emph{J. Sci. Comput.} 68 (2016) 975--1001.
\bibitem{CarpenterKennedy1994} M.\,H.~Carpenter, C.\,A.~Kennedy, Fourth-order 2N-storage Runge--Kutta schemes, NASA
 TM-109112, 1994.
\bibitem{CelledoniMarthinsenOwren2003} E.~Celledoni, A.~Marthinsen, B.~Owren, Commutator-free Lie group methods,
 \emph{Future Gener. Comput. Syst.} 19 (2003) 341--352.
\bibitem{ChinomonaReynolds2021} R.~Chinomona, D.\,R.~Reynolds, Implicit-explicit multirate infinitesimal GARK methods,
 \emph{SIAM J. Sci. Comput.} 43 (2021) A3082--A3113.
\bibitem{CoxMatthews2002} S.\,M.~Cox, P.\,C.~Matthews, Exponential time differencing for stiff systems, \emph{J.
 Comput. Phys.} 176 (2002) 430--455.
\bibitem{GiraldoKellyConstantinescu2013} F.\,X.~Giraldo, J.\,F.~Kelly, E.\,M.~Constantinescu, Implicit-explicit
 formulations of a three-dimensional nonhydrostatic unified model of the atmosphere (NUMA), \emph{SIAM J. Sci. Comput.}
 35 (2013) B1162--B1194.
\bibitem{GuntherSandu2016} M.~G\"unther, A.~Sandu, Multirate generalized additive Runge--Kutta methods, \emph{Numer.
 Math.} 133 (2016) 497--524.
\bibitem{HairerNorsettWanner1993} E.~Hairer, S.\,P.~N{\o}rsett, G.~Wanner, \emph{Solving Ordinary Differential
 Equations I}, 2nd ed., Springer, 1993.
\bibitem{HochbruckOstermann2005} M.~Hochbruck, A.~Ostermann, Explicit exponential Runge--Kutta methods for semilinear
 parabolic problems, \emph{SIAM J. Numer. Anal.} 43 (2005) 1069--1090.
\bibitem{HochbruckOstermann2010} M.~Hochbruck, A.~Ostermann, Exponential integrators, \emph{Acta Numer.} 19 (2010)
 209--286.
\bibitem{HochbruckOstermannSchweitzer2009} M.~Hochbruck, A.~Ostermann, J.~Schweitzer, Exponential Rosenbrock-type
 methods, \emph{SIAM J. Numer. Anal.} 47 (2009) 786--803.
\bibitem{KetchesonRanocha2023} D.\,I.~Ketcheson, H.~Ranocha, Computing with B-series, \emph{ACM Trans. Math. Softw.} 49
 (2023) 14.
\bibitem{KlempSkamarockDudhia2007} J.\,B.~Klemp, W.\,C.~Skamarock, J.~Dudhia, Conservative split-explicit time
 integration methods for the compressible nonhydrostatic equations, \emph{Mon. Wea. Rev.} 135 (2007) 2897--2913.
\bibitem{KlempWilhelmson1978} J.\,B.~Klemp, R.\,B.~Wilhelmson, The simulation of three-dimensional convective storm
 dynamics, \emph{J. Atmos. Sci.} 35 (1978) 1070--1096.
\bibitem{KnothWensch2014} O.~Knoth, J.~Wensch, Generalized split-explicit Runge--Kutta methods for the compressible
 Euler equations, \emph{Mon. Wea. Rev.} 142 (2014) 2067--2081.
\bibitem{KnothWolke1998} O.~Knoth, R.~Wolke, Implicit-explicit Runge--Kutta methods for computing atmospheric reactive
 flows, \emph{Appl. Numer. Math.} 28 (1998) 327--341.
\bibitem{Krogstad2005} S.~Krogstad, Generalized integrating factor methods for stiff PDEs, \emph{J. Comput. Phys.} 203
 (2005) 72--88.
\bibitem{Lawson1967} J.\,D.~Lawson, Generalized Runge--Kutta processes for stable systems with large Lipschitz
 constants, \emph{SIAM J. Numer. Anal.} 4 (1967) 372--380.
\bibitem{Luescher2010} M.~L\"uscher, Properties and uses of the Wilson flow in lattice QCD, \emph{J. High Energy
 Phys.} 08 (2010) 071.
\bibitem{LuanChinomonaReynolds2020} V.\,T.~Luan, R.~Chinomona, D.\,R.~Reynolds, A new class of high-order methods for
 multirate differential equations, \emph{SIAM J. Sci. Comput.} 42 (2020) A1245--A1268.
\bibitem{Owren2006} B.~Owren, Order conditions for commutator-free Lie group methods, \emph{J. Phys. A} 39 (2006)
 5585--5599.
\bibitem{Sandu2019} A.~Sandu, A class of multirate infinitesimal GARK methods, \emph{SIAM J. Numer. Anal.} 57 (2019)
 2300--2327.
\bibitem{WenschKnothGalant2009} J.~Wensch, O.~Knoth, A.~Galant, Multirate infinitesimal step methods for atmospheric
 flow simulation, \emph{BIT} 49 (2009) 449--473.
\bibitem{Williamson1980} J.\,H.~Williamson, Low-storage Runge--Kutta schemes, \emph{J. Comput. Phys.} 35 (1980) 48--56.
\bibitem{WickerSkamarock2002} L.\,J.~Wicker, W.\,C.~Skamarock, Time-splitting methods for elastic models using forward
 time schemes, \emph{Mon. Wea. Rev.} 130 (2002) 2088--2097.
\end{thebibliography}

\end{document}
